Let \(\mathcal O\) be the finite support of \(O\), put \(J=|\mathcal O|\), and write
\[
v=\bigl(P_{\mathrm{obs}}(O=o)\bigr)_{o\in\mathcal O}\in\mathbb R^J,
\qquad
\widehat v_n=\bigl(\widehat P_n(O=o)\bigr)_{o\in\mathcal O}.
\]
Assumption \(\mathrm{ass{:}iid\text{-}sampling}\) and Lemma \(\mathrm{lem{:}finite\text{-}multinomial\text{-}central\text{-}limit}\) give
\[
\sqrt n(\widehat v_n-v)\rightsquigarrow\mathbb Z_{P_{\mathrm{obs}}}
\quad\text{in }\mathbb R^J.
\tag{1}
\]

We first construct the reduced-support map to which the directional delta method will be applied.  Set \(\mathbb D=\mathbb R^J\), \(\mathbb E=\mathbb R^2\), and \(\mathbb D_0=\mathbb D\).  For each \(w\in\mathbb D\), compute from its coordinates the empirical-ratio analogues prescribed by Definition \(\mathrm{def{:}plugin\text{-}endpoint\text{-}estimator}\), using zero when a conditioning denominator is zero.  Form the raw capacity contrasts, apply the coordinatewise projection \(\Pi_{\mathfrak C}\), and then compute \(q_0,q_1,\Delta q,m,B_L,B_U\).  Sum numerator and denominator contributions only over the fixed set \(\mathcal X_+(P)\).  Divide the two numerators by their common denominator when it is nonzero and use the fixed value \((0,1)\) otherwise.  Denote this everywhere-defined map by
\[
\Phi:\mathbb D_\phi=\mathbb D\longrightarrow\mathbb E.
\tag{2}
\]
The arbitrary values used at zero denominators only totalize \(\Phi\) away from the population neighborhood; they are not additional probabilistic assumptions.

If \(x\in\mathcal X_+(P)\), then \(p_xm(x)>0\), and therefore \(p_x>0\).  Assumption \(\mathrm{ass{:}instrument\text{-}overlap}\) gives
\[
P_{\mathrm{obs}}(X=x,Z=1)=p_x\pi(x)\ge p_x\varepsilon_Z>0,
\qquad
P_{\mathrm{obs}}(X=x,Z=0)=p_x\{1-\pi(x)\}\ge p_x\varepsilon_Z>0.
\tag{3}
\]
All conditional-probability denominators used by \(\Phi\) on this fixed support are thus positive at \(v\) and remain nonzero in a neighborhood of \(v\).  On that neighborhood, the map from \(w\) to the finite vector \((p_x,\widetilde\ell(x),\widetilde h(x))_{x\in\mathcal X_+(P)}\) is continuously differentiable, because its coordinates are linear maps and ratios whose denominators are nonzero.

At the population law, the raw capacity vector is the nonnegative vector \(c=(\ell,h)\).  For a direction \(g\) in raw-capacity coordinates, the coordinatewise positive-part projection has Hadamard directional derivative
\[
[\Pi'_{\mathfrak C,c}(g)]_k=
\begin{cases}
g_k,&c_k>0,\\
\max\{g_k,0\},&c_k=0.
\end{cases}
\tag{4}
\]
Indeed, if \(t_r\downarrow0\) and \(g_r\to g\), then
\[
\frac{\max\{c_k+t_rg_{r,k},0\}-c_k}{t_r}
\longrightarrow
\begin{cases}
g_k,&c_k>0,\\
\max\{g_k,0\},&c_k=0.
\end{cases}
\]
There are finitely many coordinates, so the coordinatewise limits establish (4) in Euclidean norm.

Let \(\dot\ell_i(x)\) and \(\dot h_j(x)\) be the capacity-direction coordinates after (4), and define
\[
\dot q_0(x)=\sum_i\dot\ell_i(x),
\qquad
\dot q_1(x)=\sum_j\dot h_j(x),
\qquad
\dot{\Delta q}(x)=\dot q_1(x)-\dot q_0(x).
\tag{5}
\]
For the lower endpoint, set
\[
A_{x,t}=\ell_{\le t}(x)-h_{\le t}(x),
\qquad
G_x=\max_{t\in\mathcal Y}A_{x,t},
\qquad
s_x=\min\{\Delta q(x),0\}.
\tag{6}
\]
The active threshold set
\[
\mathcal A_L(x)=\{t\in\mathcal Y:A_{x,t}=G_x\}
\]
is nonempty because \(\mathcal Y\) is finite.  With
\[
\dot A_{x,t}=\sum_{i\le t}\dot\ell_i(x)-\sum_{j\le t}\dot h_j(x),
\]
the finite active-set expansion gives
\[
\dot G_x=\max_{t\in\mathcal A_L(x)}\dot A_{x,t}.
\tag{7}
\]
To verify (7), consider any \(t_r\downarrow0\) and convergent sequence of directions.  Every inactive threshold has a strictly positive gap below \(G_x\); finiteness makes the minimum inactive gap positive, so no inactive threshold can maximize for all sufficiently large \(r\).  Among the active thresholds, division by \(t_r\) leaves the maximum of the convergent first-order terms, which is (7).  Direct scalar algebra likewise gives
\[
\dot s_x=
\begin{cases}
\dot{\Delta q}(x),&\Delta q(x)<0,\\
0,&\Delta q(x)>0,\\
\min\{\dot{\Delta q}(x),0\},&\Delta q(x)=0.
\end{cases}
\tag{8}
\]
By Definition \(\mathrm{def{:}threshold\text{-}cuts}\), \(B_L(x)=\max\{0,G_x+s_x\}\), and hence
\[
\dot B_L(x)=
\begin{cases}
\dot G_x+\dot s_x,&G_x+s_x>0,\\
0,&G_x+s_x<0,\\
\max\{\dot G_x+\dot s_x,0\},&G_x+s_x=0.
\end{cases}
\tag{9}
\]
Equations (7)--(9) allow every simultaneous threshold tie.  In particular, at \(\Delta q(x)=0\), equation (8) is exactly the asserted active-face map \(g\mapsto\min\{\dot{\Delta q}(x),0\}\).

For the upper endpoint, define
\[
U_{x,t}=\ell_{<t}(x)+h_{>t}(x),
\qquad
\dot U_{x,t}=\sum_{i<t}\dot\ell_i(x)+\sum_{j>t}\dot h_j(x).
\tag{10}
\]
The survivor-mass derivative is
\[
\dot m(x)=
\begin{cases}
\dot q_0(x),&q_0(x)<q_1(x),\\
\dot q_1(x),&q_1(x)<q_0(x),\\
\min\{\dot q_0(x),\dot q_1(x)\},&q_0(x)=q_1(x).
\end{cases}
\tag{11}
\]
Index the mass cap by \(\star\), put \(V_{x,\star}=m(x)\) and \(\dot V_{x,\star}=\dot m(x)\), and put \(V_{x,t}=U_{x,t}\) and \(\dot V_{x,t}=\dot U_{x,t}\) for \(t\in\mathcal Y\).  The active upper-cut set is
\[
\mathcal A_U(x)=\{a\in\{\star\}\cup\mathcal Y:V_{x,a}=B_U(x)\}.
\]
The same finite-gap argument, applied to a minimum, yields
\[
\dot B_U(x)=\min_{a\in\mathcal A_U(x)}\dot V_{x,a}.
\tag{12}
\]
Thus (11)--(12) include positive-survivor zero-gap cells as well as arbitrary ties between the mass cap and multiple threshold cuts.

The preceding sequential expansions prove Hadamard directional differentiability rather than merely pointwise directional differentiability: each formula holds along every \(t_r\downarrow0\) and every convergent sequence of directions.  Their derivatives are finite maxima or minima of continuous linear maps and coordinatewise positive parts, and therefore are continuous functions of the direction.  Finite sums preserve these properties.  If \(\dot p_x\) denotes the cell-probability direction, then, for \(e\in\{L,U\}\), define
\[
N_e=\sum_{x\in\mathcal X_+(P)}p_xB_e(x),
\qquad
\dot N_e=\sum_{x\in\mathcal X_+(P)}
\{\dot p_xB_e(x)+p_x\dot B_e(x)\},
\tag{13}
\]
and
\[
M=\sum_{x\in\mathcal X_+(P)}p_xm(x),
\qquad
\dot M=\sum_{x\in\mathcal X_+(P)}
\{\dot p_xm(x)+p_x\dot m(x)\}.
\tag{14}
\]
Cells outside \(\mathcal X_+(P)\) have \(p_xm(x)=0\), so the first quantity in (14) equals the aggregate survivor-complier mass defined in the frozen core.  Assumption \(\mathrm{ass{:}positive\text{-}aggregate\text{-}survivors}\) gives \(M>0\), and hence the denominator remains nonzero in a neighborhood of \(v\).  The quotient expansion is therefore
\[
\left(\frac{N_e}{M}\right)'(g)
=\frac{\dot N_eM-N_e\dot M}{M^2},
\qquad e\in\{L,U\}.
\tag{15}
\]
Equations (3)--(15) establish that \(\Phi\) is Hadamard directionally differentiable at \(v\), tangentially to \(\mathbb D_0=\mathbb D\), and give its continuous derivative explicitly.

We now verify every premise of Lemma \(\mathrm{lem{:}hadamard\text{-}directional\text{-}delta\text{-}method}\).  The spaces \(\mathbb D=\mathbb R^J\) and \(\mathbb E=\mathbb R^2\) are Banach spaces; \(\mathbb D_\phi=\mathbb D\); \(\theta_0=v\); \(\widehat\theta_n=\widehat v_n\); and \(r_n=\sqrt n\uparrow\infty\).  The estimators are \(\mathbb D_\phi\)-valued.  Weak convergence is (1).  The finite-dimensional Gaussian vector \(\mathbb Z_{P_{\mathrm{obs}}}\) is tight, and its support lies in \(\mathbb R^J=\mathbb D_0\).  The required Hadamard directional differentiability and continuity of the derivative were proved above.  The cited lemma consequently gives
\[
\sqrt n\{\Phi(\widehat v_n)-\Phi(v)\}
=\Phi'_v\!\left(\sqrt n(\widehat v_n-v)\right)+o_P(1)
\rightsquigarrow\Phi'_v(\mathbb Z_{P_{\mathrm{obs}}}).
\tag{16}
\]
Because \(\mathbb D_0=\mathbb D\), every finite-sample direction already lies in the tangent set; no continuous-extension convention is needed in this application.

It remains to compare the fixed-support map (2) with the screened estimator.  Write \(r_x=p_xm(x)\).  Assumption \(M>0\) makes \(\mathcal X_+(P)\) nonempty, and finiteness of \(\mathcal X\) gives
\[
\rho=\min_{x\in\mathcal X_+(P)}r_x>0.
\tag{17}
\]
Weak convergence in (1) implies tightness of \(\sqrt n(\widehat v_n-v)\), and hence \(\widehat v_n\to_Pv\).  For every \(x\in\mathcal X_+(P)\), equation (3), continuity of the local ratio maps, and continuity of coordinatewise positive part and finite minima therefore imply \(\widehat r_x\to_Pr_x\).  Since \(\eta_n\downarrow0\), equation (17) and a finite union bound show that all these cells are retained with probability tending to one.

If \(p_x=0\), then independent sampling gives \(\widehat p_x=0\) almost surely, and hence \(\widehat r_x=0\).  Now suppose \(p_x>0\) and \(m(x)=0\).  Because \(q_0(x),q_1(x)\ge0\), at least one total, denoted by \(q_a(x)\), is zero.  Its constituent population capacities are nonnegative and have zero sum, so every one of them is zero.  Overlap and \(p_x>0\) give positive population conditioning denominators.  The corresponding raw empirical contrasts are therefore smooth functions of \(\widehat v_n\) near \(v\), centered at zero; equation (1) makes each of them \(O_P(n^{-1/2})\).  Coordinatewise positive part preserves this order, and \(K\) is fixed.  Consequently,
\[
0\le\widehat m(x)\le\widehat q_a(x)=O_P(n^{-1/2}),
\qquad
0\le\widehat r_x\le\widehat m(x)=O_P(n^{-1/2})=o_P(\eta_n),
\tag{18}
\]
where the last equality uses \(\sqrt n\eta_n\to\infty\).  A finite union over all cells combines (17)--(18) to prove
\[
P\!\left\{\{x:\widehat r_x>\eta_n\}=\mathcal X_+(P)\right\}\longrightarrow1.
\tag{19}
\]

On the event in (19), the three screened sums from Definition \(\mathrm{def{:}plugin\text{-}endpoint\text{-}estimator}\) are exactly the fixed-support sums defining \(\Phi(\widehat v_n)\).  Their common denominator is positive because the retained set is nonempty and every retained \(\widehat r_x\) is greater than \(\eta_n>0\); the fallback is therefore not used.  At the population law, Definitions \(\mathrm{def{:}threshold\text{-}cuts}\) and \(\mathrm{def{:}identified\text{-}interval}\), together with Theorem \(\mathrm{thm{:}sharp\text{-}exact\text{-}mass\text{-}threshold\text{-}interval}\), give \(\Phi(v)=\Psi(P_{\mathrm{obs}})\).  Replacing \(\Phi(\widehat v_n)\) by \(\widehat\Psi_n\) in (16) on an event whose probability tends to one proves the stated directional limit.  Continuity of \(\Phi\) at \(v\), the consequence \(\widehat v_n\to_Pv\) of (1), and the same replacement prove \(\widehat\Psi_n\to_P\Psi(P_{\mathrm{obs}})\).  No direction-separation, multiplier-calibration, or uniform-in-law conclusion is used. \(\square\)